% Figura: hmi-datos -- requiere % =====================================================================
%  Estilos compartidos por los diagramas TikZ de los capitulos.
%  Va en el PREAMBULO del documento principal:
%      % =====================================================================
%  Estilos compartidos por los diagramas TikZ de los capitulos.
%  Va en el PREAMBULO del documento principal:
%      % =====================================================================
%  Estilos compartidos por los diagramas TikZ de los capitulos.
%  Va en el PREAMBULO del documento principal:
%      \input{diagramas/estilos-diagramas}
% =====================================================================

\usepackage{tikz}
\usepackage{amsmath}
\usepackage{amssymb}
\usetikzlibrary{arrows.meta,positioning,shapes.geometric,shapes.misc,calc,fit,backgrounds,decorations.pathreplacing}

\definecolor{cbloque}{HTML}{E8EEF4}
\definecolor{cborde}{HTML}{4A6D8C}
\definecolor{cacento}{HTML}{D9E6D5}
\definecolor{cacentob}{HTML}{5A7A4A}
\definecolor{cgris}{HTML}{EDEDED}
\definecolor{cgrisb}{HTML}{888888}
\definecolor{calarma}{HTML}{F6DEDE}
\definecolor{calarmab}{HTML}{A03030}

\tikzset{
  bloque/.style={draw=cborde,fill=cbloque,rounded corners=2pt,align=center,
                 inner sep=5pt,font=\footnotesize,minimum height=9mm},
  bloqueg/.style={bloque,draw=cgrisb,fill=cgris},
  bloquea/.style={bloque,draw=cacentob,fill=cacento},
  bloquer/.style={bloque,draw=calarmab,fill=calarma},
  proceso/.style={bloque,minimum width=52mm},
  decision/.style={draw=cborde,fill=cbloque,align=center,font=\scriptsize,
                   diamond,aspect=2.2,inner sep=1pt},
  flecha/.style={-{Stealth[length=2mm]},draw=cborde!80!black,thick},
  flechag/.style={flecha,draw=cgrisb},
  et/.style={font=\scriptsize,inner sep=1.5pt,fill=white,align=center},
  etl/.style={et,midway},
  grupo/.style={draw=cgrisb,dashed,rounded corners=3pt,inner sep=4mm},
  tit/.style={font=\scriptsize\bfseries,text=cgrisb!60!black},
}

% --- utilidades para diagramas de secuencia ---
\tikzset{
  vida/.style={draw=cgrisb,dash pattern=on 2pt off 2pt},
  cab/.style={bloque,minimum width=26mm,font=\scriptsize\bfseries},
  nota/.style={draw=cacentob,fill=cacento,font=\scriptsize,align=center,
               rounded corners=1pt,inner sep=3pt},
}
\newcommand{\msg}[4]{\draw[flecha] (#1,#3) -- (#2,#3) node[et,midway,above=0pt]{#4};}
\newcommand{\msgr}[4]{\draw[flecha,dashed] (#1,#3) -- (#2,#3) node[et,midway,above=0pt]{#4};}
\newcommand{\auto}[3]{%
  \draw[flecha] (#1,#2) -- ++(0.55,0) -- ++(0,-0.42) -- (#1,{#2-0.42});
  \node[et,anchor=west] at ({#1+0.7},{#2-0.21}) {#3};}


% =====================================================================

\usepackage{tikz}
\usepackage{amsmath}
\usepackage{amssymb}
\usetikzlibrary{arrows.meta,positioning,shapes.geometric,shapes.misc,calc,fit,backgrounds,decorations.pathreplacing}

\definecolor{cbloque}{HTML}{E8EEF4}
\definecolor{cborde}{HTML}{4A6D8C}
\definecolor{cacento}{HTML}{D9E6D5}
\definecolor{cacentob}{HTML}{5A7A4A}
\definecolor{cgris}{HTML}{EDEDED}
\definecolor{cgrisb}{HTML}{888888}
\definecolor{calarma}{HTML}{F6DEDE}
\definecolor{calarmab}{HTML}{A03030}

\tikzset{
  bloque/.style={draw=cborde,fill=cbloque,rounded corners=2pt,align=center,
                 inner sep=5pt,font=\footnotesize,minimum height=9mm},
  bloqueg/.style={bloque,draw=cgrisb,fill=cgris},
  bloquea/.style={bloque,draw=cacentob,fill=cacento},
  bloquer/.style={bloque,draw=calarmab,fill=calarma},
  proceso/.style={bloque,minimum width=52mm},
  decision/.style={draw=cborde,fill=cbloque,align=center,font=\scriptsize,
                   diamond,aspect=2.2,inner sep=1pt},
  flecha/.style={-{Stealth[length=2mm]},draw=cborde!80!black,thick},
  flechag/.style={flecha,draw=cgrisb},
  et/.style={font=\scriptsize,inner sep=1.5pt,fill=white,align=center},
  etl/.style={et,midway},
  grupo/.style={draw=cgrisb,dashed,rounded corners=3pt,inner sep=4mm},
  tit/.style={font=\scriptsize\bfseries,text=cgrisb!60!black},
}

% --- utilidades para diagramas de secuencia ---
\tikzset{
  vida/.style={draw=cgrisb,dash pattern=on 2pt off 2pt},
  cab/.style={bloque,minimum width=26mm,font=\scriptsize\bfseries},
  nota/.style={draw=cacentob,fill=cacento,font=\scriptsize,align=center,
               rounded corners=1pt,inner sep=3pt},
}
\newcommand{\msg}[4]{\draw[flecha] (#1,#3) -- (#2,#3) node[et,midway,above=0pt]{#4};}
\newcommand{\msgr}[4]{\draw[flecha,dashed] (#1,#3) -- (#2,#3) node[et,midway,above=0pt]{#4};}
\newcommand{\auto}[3]{%
  \draw[flecha] (#1,#2) -- ++(0.55,0) -- ++(0,-0.42) -- (#1,{#2-0.42});
  \node[et,anchor=west] at ({#1+0.7},{#2-0.21}) {#3};}


% =====================================================================

\usepackage{tikz}
\usepackage{amsmath}
\usepackage{amssymb}
\usetikzlibrary{arrows.meta,positioning,shapes.geometric,shapes.misc,calc,fit,backgrounds,decorations.pathreplacing}

\definecolor{cbloque}{HTML}{E8EEF4}
\definecolor{cborde}{HTML}{4A6D8C}
\definecolor{cacento}{HTML}{D9E6D5}
\definecolor{cacentob}{HTML}{5A7A4A}
\definecolor{cgris}{HTML}{EDEDED}
\definecolor{cgrisb}{HTML}{888888}
\definecolor{calarma}{HTML}{F6DEDE}
\definecolor{calarmab}{HTML}{A03030}

\tikzset{
  bloque/.style={draw=cborde,fill=cbloque,rounded corners=2pt,align=center,
                 inner sep=5pt,font=\footnotesize,minimum height=9mm},
  bloqueg/.style={bloque,draw=cgrisb,fill=cgris},
  bloquea/.style={bloque,draw=cacentob,fill=cacento},
  bloquer/.style={bloque,draw=calarmab,fill=calarma},
  proceso/.style={bloque,minimum width=52mm},
  decision/.style={draw=cborde,fill=cbloque,align=center,font=\scriptsize,
                   diamond,aspect=2.2,inner sep=1pt},
  flecha/.style={-{Stealth[length=2mm]},draw=cborde!80!black,thick},
  flechag/.style={flecha,draw=cgrisb},
  et/.style={font=\scriptsize,inner sep=1.5pt,fill=white,align=center},
  etl/.style={et,midway},
  grupo/.style={draw=cgrisb,dashed,rounded corners=3pt,inner sep=4mm},
  tit/.style={font=\scriptsize\bfseries,text=cgrisb!60!black},
}

% --- utilidades para diagramas de secuencia ---
\tikzset{
  vida/.style={draw=cgrisb,dash pattern=on 2pt off 2pt},
  cab/.style={bloque,minimum width=26mm,font=\scriptsize\bfseries},
  nota/.style={draw=cacentob,fill=cacento,font=\scriptsize,align=center,
               rounded corners=1pt,inner sep=3pt},
}
\newcommand{\msg}[4]{\draw[flecha] (#1,#3) -- (#2,#3) node[et,midway,above=0pt]{#4};}
\newcommand{\msgr}[4]{\draw[flecha,dashed] (#1,#3) -- (#2,#3) node[et,midway,above=0pt]{#4};}
\newcommand{\auto}[3]{%
  \draw[flecha] (#1,#2) -- ++(0.55,0) -- ++(0,-0.42) -- (#1,{#2-0.42});
  \node[et,anchor=west] at ({#1+0.7},{#2-0.21}) {#3};}

 en el preambulo
\begin{tikzpicture}[x=1cm,y=1cm,
  paso/.style={bloque,text width=54mm},
  dec/.style={decision,text width=30mm},
  canal/.style={draw=cborde!80!black,thick},
  banda/.style={draw=cgrisb,dashed,rounded corners=3pt,fill=cgris!45}]

% ---------------- bandas ----------------
\begin{scope}[on background layer]
  \fill[banda] (-9.6,-9.4) rectangle (11.0,2.9);
  \fill[banda] (-9.6,-35.2) rectangle (11.0,-12.4);
\end{scope}
\node[tit,anchor=north west] at (-9.5,2.8)   {hilo lector \textemdash\ \texttt{serie.py}, continuo};
\node[tit,anchor=north west] at (-9.5,-12.5) {hilo principal de Tk \textemdash\ \texttt{hmi.py}};

% ---------------- hilo lector ----------------
\node[dec]  (a1) at (0,0)    {¿hay bytes\\en el puerto?};
\node[paso] (a2) at (0,-2.4) {Leer, decodificar y acumular en el buffer};
\node[dec]  (a3) at (0,-4.8) {¿el buffer tiene\\una línea completa?};
\node[paso] (a4) at (0,-7.4) {Cortar la línea y ponerla en la cola};

\node[bloqueg,text width=30mm] (aw) at (6.4,0)  {Esperar 5\,ms};
\node[bloquer,text width=34mm] (ae) at (-6.6,-2.4)
  {Falla del puerto:\\[1pt]\scriptsize levantar \texttt{connection\_lost} y terminar el hilo};

\draw[flecha] (a1) -- node[et,right]{sí} (a2);
\draw[flecha] (a2) -- (a3);
\draw[flecha] (a3) -- node[et,right]{sí} (a4);
\draw[flecha] (a1.east) -- node[et,above]{no} (aw.west);
\draw[canal] (aw.north) -- (6.4,1.6) -- (0,1.6);
\draw[canal] (a3.west) -- node[et,above,pos=0.25]{no} (-9.1,-4.8) -- (-9.1,1.6) -- (0,1.6);
\draw[flecha] (0,1.6) -- (a1.north);
\draw[flecha] (a4.east) -- (9.6,-7.4) -- (9.6,-4.8) -- (a3.east);
\draw[flechag] (a2.west) -- node[et,above]{excepción} (ae.east);

% ---------------- la cola ----------------
\node[bloquea,text width=58mm,draw=cacentob,dashed] (q) at (0,-10.9)
  {\texttt{queue.Queue}\\[1pt]\scriptsize frontera entre los dos hilos};
\draw[flecha] (a4.south) -- (q.north);
\draw[flecha] (q.south) -- (0,-13.05);

% ---------------- hilo principal de Tk ----------------
\node[bloquea,text width=54mm] (b0) at (0,-13.8)
  {\texttt{\_poll\_serial()}\\[1pt]\scriptsize \texttt{root.after} cada 100\,ms};
\node[paso] (b1) at (0,-16.2) {Tomar de la cola hasta 50 líneas};
\node[dec]  (b2) at (0,-18.6) {¿queda alguna\\línea del lote?};
\node[paso] (b3) at (0,-21.2) {Eco en el monitor y \texttt{parse\_line()}\\[1pt]\scriptsize actualiza \texttt{ProcessData}};
\node[dec]  (b4) at (0,-23.8) {¿la línea\\cerró un ciclo?};
\node[paso,text width=58mm] (b5) at (0,-26.4)
  {\texttt{track\_minima()}, punto al gráfico, fila al CSV y paso de la sintonización automática};
\node[dec]  (b6) at (0,-29.0) {¿\texttt{connection}\\\texttt{\_lost}?};
\node[dec]  (b7) at (0,-32.2) {¿hubo\\novedades?};

\node[bloquer,text width=34mm] (bd) at (6.6,-29.0) {Avisar y cerrar el puerto};
\node[bloque,text width=34mm]  (br) at (6.6,-32.2) {\texttt{\_refresh\_display()}};

\draw[flecha] (b0) -- (b1);
\draw[flecha] (b1) -- (b2);
\draw[flecha] (b2) -- node[et,right]{sí} (b3);
\draw[flecha] (b3) -- (b4);
\draw[flecha] (b4) -- node[et,right]{sí} (b5);
\draw[canal] (b4.east) -- node[et,above,pos=0.35]{no} (9.6,-23.8);
\draw[canal] (b5.east) -- (9.6,-26.4);
\draw[canal] (9.6,-26.4) -- (9.6,-18.6);
\draw[flecha] (9.6,-18.6) -- (b2.east);
\draw[flecha] (b2.west) -- node[et,above,pos=0.22]{no} (-9.1,-18.6) -- (-9.1,-29.0) -- (b6.west);
\draw[flecha] (b6) -- node[et,left]{no} (b7);
\draw[flecha] (b6.east) -- node[et,above]{sí} (bd.west);
\draw[flecha] (b7.east) -- node[et,above]{sí} (br.west);
\draw[flechag] (bd.south) -- (6.6,-30.6) -- (0,-30.6);
\draw[flechag] (br.south) -- (6.6,-34.2) -- (0,-34.2);
\draw[flecha] (b7.south) -- node[et,left,pos=0.45]{no} (0,-34.2) -- (10.3,-34.2)
      -- node[et,pos=0.55,right,rotate=90]{\texttt{root.after(100\,ms)}} (10.3,-13.8) -- (b0.east);

\end{tikzpicture}